% Part: incompleteness
% Chapter: representability-in-q
% Section: comp-representable

\documentclass[../../../include/open-logic-section]{subfiles}

\begin{document}

\olfileid[hi]{inc}{req}{crq}
\olsection{संगणनीय फलन $\Th{Q}$ में निरूपणीय हैं}

\begin{thm}
प्रत्येक संगणनीय फलन~$\Th{Q}$ में निरूपणीय है।
\end{thm}

\begin{proof}
निश्चितता के लिए और चर्च--ट्यूरिंग अभिधारणा का उपयोग करते हुए, हम कहें कि
कोई फलन तभी और केवल तभी संगणनीय है जब वह सामान्य पुनरावर्ती है। सामान्य
पुनरावर्ती फलन वे हैं जिन्हें शून्य फलन~$\Zero$, उत्तराधिकारी फलन~$\Succ$
और प्रक्षेप फलन~$\Proj{n}{i}$ से संयोजन, आदिम पुनरावर्तन तथा नियमित
न्यूनीकरण का उपयोग करके परिभाषित किया जा सकता है।
\olref[pri]{lem:prim-rec} के अनुसार $f$ और~$g$ से परिभाषित किए जा सकने वाले
किसी भी फलन~$h$ को $f$, $g$ तथा $\Zero$, $\Succ$, $\Proj{n}{i}$, $\Add$,
$\Mult$, $\Char{=}$ से संयोजन और नियमित न्यूनीकरण द्वारा भी परिभाषित किया
जा सकता है। फलतः कोई फलन सामान्य पुनरावर्ती तभी और केवल तभी है जब उसे
$\Zero$, $\Succ$, $\Proj{n}{i}$, $\Add$, $\Mult$, $\Char{=}$ से संयोजन तथा
नियमित न्यूनीकरण द्वारा परिभाषित किया जा सके।

इसके अतिरिक्त हमने दिखाया है कि विचाराधीन मूल फलन~$\Th{Q}$ में निरूपणीय
हैं
(\cref{inc:req:bre:prop:rep-zero,inc:req:bre:prop:rep-succ,inc:req:bre:prop:rep-proj,inc:req:bre:prop:rep-id,inc:req:bre:prop:rep-add,inc:req:bre:prop:rep-mult}),
और निरूपणीय फलनों से संयोजन अथवा नियमित न्यूनीकरण द्वारा परिभाषित कोई भी
फलन
(\olref[cmp]{prop:rep-composition},
\olref[min]{prop:rep-minimization}) भी निरूपणीय है। अतः प्रत्येक सामान्य
पुनरावर्ती फलन~$\Th{Q}$ में निरूपणीय है।
\end{proof}

\begin{explain}
हमने दिखाया है कि संगणनीय फलनों के समुच्चय को $\Th{Q}$ में निरूपणीय फलनों
के समुच्चय के रूप में अभिलक्षित किया जा सकता है। वास्तव में प्रमाण इससे
अधिक व्यापक है। निरूपणीयता की परिभाषा से यह देखना कठिन नहीं है कि
$\Th{Q}$ का कोई भी विस्तार (अथवा कोई ऐसा सिद्धांत जिसमें $\Th{Q}$ की
व्याख्या की जा सके) संगणनीय फलनों का निरूपण कर सकता है। किंतु विलोमतः,
जिस किसी !!{derivation}-तंत्र में !!{derivation} की धारणा संगणनीय हो, उसमें
प्रत्येक निरूपणीय फलन संगणनीय है। अतः, उदाहरण के लिए, संगणनीय फलनों के
समुच्चय को पियानो अंकगणित में, बल्कि ज़र्मेलो--फ्रेंकेल समुच्चय-सिद्धांत
में भी, निरूपणीय फलनों के समुच्चय के रूप में अभिलक्षित किया जा सकता है।
जैसा ग्योडेल (G\"odel) ने उल्लेख किया, यह कुछ आश्चर्यजनक है। हम देखेंगे कि
सिद्धता के विषय में प्रश्न इस बात के प्रति अत्यंत संवेदनशील होते हैं कि हम
किस सिद्धांत पर विचार कर रहे हैं; मोटे तौर पर, अभिगृहीत जितने प्रबल होंगे,
उतना अधिक सिद्ध किया जा सकेगा। किंतु अभिगृहीती सिद्धांतों के एक व्यापक
वर्ग में निरूपणीय फलन ठीक संगणनीय फलन ही हैं; जब तक सिद्धांत
अभिगृहीतीकरणीय हैं, अधिक प्रबल सिद्धांत अधिक फलनों का निरूपण नहीं करते।
\end{explain}

\end{document}
